% Section 2 — the algebra of false agreement.
%
% Scope note, to keep the paper honest: the identity below is elementary and the
% correlation floor is standard ensemble theory (bias--variance--covariance).
% Neither is claimed as a contribution. What is new here is (i) applying them to
% probe-based cross-model oversight, (ii) the transport term of
% Sec.~\ref{sec:theory:transport}, and (iii) measuring the error correlation
% directly on internal readouts rather than on outputs. The section is framing
% that motivates a measurement, not a theorem.
%
% Every quantity in this section has an implementation in cmb/metrics.py; the
% margin comments name it.
\section{The Algebra of False Agreement}
\label{sec:theory}

\subsection{Setup}
\label{sec:theory:setup}

Restrict attention to the statements whose ground truth is \textsc{false}; these
are the only items on which false agreement is defined. For probe $i \in \{1,2\}$
let
\[
  E_i = \mathbf{1}\{\text{probe $i$ reads \textsc{true}}\},
  \qquad
  p_i = \Prob(E_i = 1),
\]
so $p_i$ is probe $i$'s false-positive rate on this restricted population.
\emph{False agreement} is the event $\{E_1 = E_2 = 1\}$, and its rate is
\[
  \FA = \Prob(E_1 = 1,\, E_2 = 1).
\]
In the eight-cell joint table of two verdicts against ground truth
(Section~\ref{sec:method}), $\FA$ is the mass of the single cell in which both
probes say \textsc{true} and the claim is false --- the one cell where
disagreement-based routing stays silent and is wrong.

\subsection{The false-agreement identity}
\label{sec:theory:identity}

\begin{proposition}[False agreement is set by error correlation]
\label{prop:identity}
Let $\rho$ be the Pearson correlation between $E_1$ and $E_2$, which for binary
variables is the $\phi$ coefficient. Then
\begin{equation}
  \FA = p_1 p_2 + \rho \sqrt{p_1(1-p_1)}\,\sqrt{p_2(1-p_2)}.
  \label{eq:identity}
\end{equation}
\end{proposition}

\begin{proof}
For indicator variables, $\mathbb{E}[E_i] = p_i$, $\mathbb{E}[E_1 E_2] =
\Prob(E_1=1, E_2=1) = \FA$, and $\mathrm{Var}(E_i) = p_i(1-p_i)$. Hence
$\Cov(E_1, E_2) = \FA - p_1 p_2$, and substituting
$\Cov(E_1,E_2) = \rho \sqrt{\mathrm{Var}(E_1)\mathrm{Var}(E_2)}$ gives
\eqref{eq:identity}.
\end{proof}

Equation~\eqref{eq:identity} is an identity, not an approximation or a bound: it
holds for every joint distribution of the two verdicts. Its content is the split
between the two terms. The first, $p_1 p_2$, is what accuracy buys. The second is
what \emph{diversity} buys, and it can dominate. Two accurate probes whose errors
are strongly correlated can suffer a higher $\FA$ than two less accurate probes
whose errors are independent --- which is why selecting overseers on accuracy
alone is not sound.

\subsection{Reference points and the feasible range of $\rho$}
\label{sec:theory:bounds}

Three reference points, stated with the running example $p_1 = 10\%$,
$p_2 = 5\%$:

\begin{description}
  \item[Independence ($\rho = 0$):] $\FA = p_1 p_2 = 0.5\%$. The errors coincide
    only by chance: of $1{,}000$ false statements, probe 2 fails on $50$, and
    probe 1 fails on $10\%$ of those, which is $5$. This is a \emph{reference
    point}, not a bound --- nothing forces two models trained on overlapping
    corpora to reach it.
  \item[Maximal overlap (upper bound):] $\FA = \min(p_1, p_2) = 5\%$. Every error
    of the better probe is also an error of the worse one.
  \item[Disjoint errors (lower bound):] $\FA = \max(0,\, p_1 + p_2 - 1) = 0$.
    This requires anti-correlated errors, which is not a realistic operating
    point for models trained on overlapping data.
\end{description}

\begin{proposition}[Fr\'echet bounds and the attainable correlation]
\label{prop:frechet}
$\max(0, p_1 + p_2 - 1) \le \FA \le \min(p_1, p_2)$, and consequently $\rho$ is
confined to $[\rho_{\min}, \rho_{\max}]$ with
\begin{equation}
  \rho_{\max} = \frac{\min(p_1,p_2) - p_1 p_2}{\sqrt{p_1(1-p_1)p_2(1-p_2)}},
  \qquad
  \rho_{\min} = \frac{\max(0,\,p_1+p_2-1) - p_1 p_2}{\sqrt{p_1(1-p_1)p_2(1-p_2)}}.
  \label{eq:rhorange}
\end{equation}
In particular $\rho$ spans the full interval $[-1, 1]$ only when $p_1 = p_2$.
\end{proposition}

\begin{proof}
The bounds on $\FA$ are the Fr\'echet--Hoeffding bounds for the joint
distribution of two indicators with fixed marginals. Inverting
\eqref{eq:identity} for $\rho$ is monotone increasing in $\FA$, so the extremes
of $\FA$ map to the extremes of $\rho$, giving \eqref{eq:rhorange}.
\end{proof}

This matters for reporting. In the running example, maximal overlap corresponds
to $\rho \approx 0.69$, not $\rho = 1$, and disjointness to $\rho \approx -0.08$.
A raw $\rho$ is therefore not comparable across pairs with different
false-positive rates. We report either $\rho$ normalized by its feasible extreme,
\begin{equation}
  \tilde\rho = \rho / \rho_{\max} \ \ (\rho \ge 0),
  \qquad
  \tilde\rho = \rho / |\rho_{\min}| \ \ (\rho < 0),
  \label{eq:rhonorm}
\end{equation}
or $\FA$ directly alongside its two Fr\'echet bounds, which is
scale-free by construction. Self-oversight --- a model checking itself --- is the
degenerate case in which $\FA$ attains its upper bound and the second reader adds
nothing.

\subsection{Detectable coverage}
\label{sec:theory:coverage}

The same $2\times 2$ table yields the positive claim and the negative one at
once. On the false statements the mass splits into \emph{agree-wrong} ($\FA$),
\emph{disagree}, and \emph{agree-right}. Disagreement is what a router can flag;
$\FA$ is the blind spot. Taking the stronger probe --- the one with the lower
false-positive rate, $p_{\mathrm{strong}} = \min(p_1,p_2)$ --- as the system
whose errors we are trying to catch:

\begin{corollary}[Coverage of a disagreement router]
\label{cor:coverage}
The fraction of the strong probe's false positives that the router flags is
\begin{equation}
  C = 1 - \frac{\FA}{p_{\mathrm{strong}}}.
  \label{eq:coverage}
\end{equation}
\end{corollary}

With the running example, $C = 90\%$ under independence and $C = 0\%$ under
maximal overlap. One experiment therefore reports both the oversight claim and
its blind spot, and they are two readings of the same table. A caveat that
survives all of this: disagreement establishes only that \emph{one} of the two
readers is wrong, never which one, so the router buys a labeled check rather than
a correction.

\subsection{Transport-adjusted false agreement}
\label{sec:theory:transport}

When the second reader is not a natively fitted probe but a probe carried across
models through the alignment map, its false-positive rate degrades:
\[
  p_2' = p_2 + \eps ,
\]
where $\eps$ is the Gate~B degradation --- native minus transported discriminative
performance, converted to an error rate at a matched threshold. We measure $\eps$
as the realized inflation of the false-positive rate at the threshold that holds
the probe's positive rate fixed, rather than reading it off an AUROC gap, so that
$\eps$ enters \eqref{eq:identity} in the same units as $p_2$.

Two consequences shape the reporting. First, $\FA$ is reported \emph{twice}, once
with native probes and once with transported ones; the gap between the two is the
cost of the map, and it belongs in the paper as a measured quantity rather than a
caveat. Second, transport plausibly inflates $\rho$ as well, because the map is
fitted from the source model's geometry and thereby imports that model's
structure into the target's readout. We therefore recompute $\rho$ on the
transported readouts instead of reusing the native $\rho$; substituting
$(p_1, p_2', \rho')$ into \eqref{eq:identity} gives the transport-adjusted rate.

\paragraph{Bidirectional intersection.}
Transport has a direction, and a false agreement that appears under only one
direction may be an artifact of the map rather than a shared belief error. We
compute the false-agreement set under both $A \to B$ and $B \to A$ transport and
treat the \emph{intersection} as the robust false-agreement set; items flagged in
one direction only are counted as transport noise, and the size of that gap (we
report the Jaccard overlap of the two sets) is itself a diagnostic of map
quality.

\subsection{Remark: the correlation floor for $N$ overseers}
\label{sec:theory:neff}

\begin{remark}[Effective number of independent overseers]
\label{rem:neff}
For $N$ overseers with equal error variance and mean pairwise error correlation
$\rhobar$, the effective number of independent overseers is
\begin{equation}
  \Neff = \frac{N}{1 + (N-1)\rhobar}
  \;\xrightarrow[N \to \infty]{}\; \frac{1}{\rhobar}.
  \label{eq:neff}
\end{equation}
\end{remark}

This is the \emph{correlation floor}: once $\rhobar$ is positive, adding
overseers stops helping, and the ceiling $1/\rhobar$ does not depend on $N$. At
$\rhobar = 0.8$ the cap is $\Neff \le 1.25$ --- adding models to a pool of
near-clones buys essentially nothing; it takes $\rhobar \approx 0.3$ for $N = 20$
readers to be worth about $3$ independent ones. The result is standard ensemble
theory (the bias--variance--covariance decomposition), and we cite it rather than
claim it. The $N$-model case is developed properly in follow-on work; here it
serves only to show that the pairwise quantity we measure, $\rho$, is the
quantity that governs the general case.
